The numerical solver in \texttt{src/DiffSol.m} approximates the model on a
uniform grid in time,
\[
t_i = t_{\min} + i h, \qquad i=0,1,\ldots,N, \qquad
h = \frac{t_{\max}-t_{\min}}{N}.
\]

After discretization, the unknown temperatures satisfy a tridiagonal linear
system. Tridiagonal systems are computationally efficient because each row has
only a small number of nonzero entries.

In this implementation, the coefficient matrix is assembled in sparse form and
solved with MATLAB's backslash operator. This combination is reliable and
efficient for the problem sizes considered in this project.

The routine returns:
\begin{itemize}
\item \texttt{timeNum}: the grid values $t_i$,
\item \texttt{TempNum}: the numerical temperature vector.
\end{itemize}
The first numerical value is fixed by the initial condition,
\texttt{TempNum(1)=T0}.
